\section*{3.5 Uniqueness of Measure Extension}

We will use the term \textit{set system} to refer to a collection of subsets in $\Omega$.

\begin{defbox}{3.10}
We define a \textit{$\pi$-system} in $\Omega$ as a set system $\mathcal{P}$ that is closed under intersection, i.e., it satisfies
\[
A \in \mathcal{P} \text{ and } B \in \mathcal{P} \Rightarrow A \cap B \in \mathcal{P}.
\]
A \textit{$\lambda$-system} in $\Omega$ is a set system $\mathcal{L}$ that satisfies the following conditions:
\begin{enumerate}[label=(\roman*)]
    \item $\Omega \in \mathcal{L}$.
    \item If $A \in \mathcal{L}$, then $\Omega \setminus A \in \mathcal{L}$.
    \item If $A_i \in \mathcal{L}$ for $i = 1, 2, 3, \dots$ is a sequence of mutually disjoint sets, then $\cup_i A_i \in \mathcal{L}$.
\end{enumerate}
\end{defbox}

A $\lambda$-system is also known as a Dynkin system, named after the mathematician Eugene Dynkin, who first introduced this concept. The $\pi$-$\lambda$ theorem, also proved by Dynkin, is a fundamental result in measure theory and probability theory.

The theorem says that if a $\pi$-system $\mathcal{P}$ is contained in a $\lambda$-system $\mathcal{L}$, then we can enlarge the $\mathcal{P}$ to the $\sigma$-field generated by $\mathcal{P}$, and this $\sigma$-field is also contained in $\mathcal{L}$. Formally, the $\pi$-$\lambda$ theorem is stated as follows.

\begin{thmbox}{3.7 ($\pi$-$\lambda$ Theorem)}
If $\mathcal{L}$ is a $\lambda$-system containing a $\pi$-system $\mathcal{P}$, then
\[
\sigma(\mathcal{P}) \subseteq \mathcal{L}.
\]
\end{thmbox}

The proof of $\pi$-$\lambda$ theorem is an exercise (Exercise 3.9).

We can use the $\pi$-$\lambda$ theorem to prove the uniqueness of measure extension.

\begin{thmbox}{3.8 (Uniqueness of Measure Extension)}
Suppose $\mathcal{F}_0$ is a field on a sample space $\Omega$, and $\mu_1$ and $\mu_2$ are two measures on $\sigma(\mathcal{F}_0)$ such that $\mu_1(A) = \mu_2(A)$ for all $A \in \mathcal{F}_0$. Furthermore, suppose there exists a sequence of disjoint sets $\Omega_i \in \mathcal{F}_0$, for $i = 1, 2, 3, \dots$, such that $\uplus_i \Omega_i = \Omega$ and $\mu_1(\Omega_i) = \mu_2(\Omega_i) < \infty$ for all $i$. Then $\mu_1(B) = \mu_2(B)$ for all $B \in \sigma(\mathcal{F}_0)$.
\end{thmbox}
The theorem states that if two measures on $\sigma(\mathcal{F}_0)$ agree on a field $\mathcal{F}_0$ that generates $\sigma(\mathcal{F}_0)$ and also agree on a sequence of disjoint sets that cover the entire sample space, then they agree on $\sigma(\mathcal{F}_0)$.

\noindent\textit{Proof} Let $S$ be a set in $\mathcal{F}_0$ such that $\mu_1(S) = \mu_2(S) < \infty$. Consider the collection of subsets
\[
\mathcal{M}_S \triangleq \{Q \in \sigma(\mathcal{F}_0) : \mu_1(S \cap Q) = \mu_2(S \cap Q)\}.
\]
We claim that
\begin{enumerate}[label=(\alph*)]
    \item $\mathcal{F}_0 \subseteq \mathcal{M}_S$, and
    \item $\mathcal{M}_S$ is a $\lambda$-system.
\end{enumerate}

The proofs of (a) and (b) will be given in the second half of the proof. We first show that the measure extension is unique given the claims in (a) and (b).

Since $\mathcal{F}_0$ is certainly a $\pi$-system, by the $\pi$-$\lambda$ theorem, we have
\begin{equation}
\mathcal{M}_S \supseteq \sigma(\mathcal{F}_0). \tag{3.3}
\end{equation}

We apply (3.3) with $S$ replaced by $\Omega_i$, for $i = 1, 2, 3, \dots$. This yields $\sigma(\mathcal{F}_0) \subseteq \mathcal{M}_{\Omega_i}$ for all $i$.

Let $E$ be any event in $\sigma(\mathcal{F}_0)$. We can write $E$ as
\[
E = E \cap \Omega = E \cap \Bigl( \biguplus_{i=1}^\infty \Omega_i \Bigr) = \biguplus_{i=1}^\infty (E \cap \Omega_i).
\]
Since $\mu_1$ and $\mu_2$ are measure functions, they enjoy the property of $\sigma$-additivity. Hence,
\[
\mu_k(E) = \sum_{i=1}^\infty \mu_k(E \cap \Omega_i)
\]
for $k = 1, 2$.

Since $E \in \sigma(\mathcal{F}_0) \subseteq \mathcal{M}_{\Omega_i}$ for all $i$, we have
\[
\mu_1(E \cap \Omega_i) = \mu_2(E \cap \Omega_i)
\]
for all $i$, by the definition of $\mathcal{M}_{\Omega_i}$. This implies that $\mu_1(E) = \mu_2(E)$ and thus completes the proof of the uniqueness of measure extension.

We now prove (a): Suppose $A \in \mathcal{F}_0$. Then $A \cap S \in \mathcal{F}_0$ because both $A$ and $S$ are in $\mathcal{F}_0$ and $\mathcal{F}_0$ is closed under intersection. By the assumption on $\mu_1$ and $\mu_2$, we have $\mu_1(A \cap S) = \mu_2(A \cap S)$. Therefore, $A \in \mathcal{M}_S$ for all $A \in \mathcal{F}_0$. This proves that $\mathcal{F}_0 \subseteq \mathcal{M}_S$.
To prove (b), we divide the argument into three parts.
\begin{enumerate}
    \item We have $\Omega \in \mathcal{M}_S$ because $\mu_1(\Omega \cap S) = \mu_1(S) = \mu_2(S) = \mu_2(\Omega \cap S)$.
    \item Suppose $A \in \mathcal{M}_S$. Because $S$ can be written as $S = (A \cap S) \uplus (A^c \cap S)$, we have $\mu_k(S) = \mu_k(A \cap S) + \mu_k(A^c \cap S)$ for $k = 1, 2$. Using the assumption that $\mu_1(S) = \mu_2(S) < \infty$, we obtain
    \[
    \mu_1(A^c \cap S) = \mu_1(S) - \mu_1(A \cap S) = \mu_2(S) - \mu_2(A \cap S) = \mu_2(A^c \cap S).
    \]
    Therefore, $A^c \in \mathcal{M}_S$.
    \item Suppose $A_i$ are mutually disjoint sets in $\mathcal{M}_S$. This means that $\mu_1(A_i \cap S) = \mu_2(A_i \cap S)$ for all $i$. By $\sigma$-additivity, we obtain
    \[
    \mu_1\Bigl( \uplus_i (A_i \cap S) \Bigr) = \mu_2\Bigl( \uplus_i (A_i \cap S) \Bigr) \implies \mu_1\Bigl( (\uplus_i A_i) \cap S \Bigr) = \mu_2\Bigl( (\uplus_i A_i) \cap S \Bigr).
    \]
    This shows that $\uplus_i A_i$ is in $\mathcal{M}_S$, completing the proof that $\mathcal{M}_S$ is a $\lambda$-system. \hfill $\blacksquare$
\end{enumerate}

As an application of the uniqueness result, we prove that a probability measure on $\mathbb{R}$ is uniquely determined by its Stieltjes measure function.

\begin{thmbox}{3.9}
Let $P$ be a probability measure on $(\mathbb{R}, \mathcal{B}(\mathbb{R}))$, and let $F(x)$ be its induced distribution function. If there is another probability $Q$ on $\mathcal{B}(\mathbb{R})$ such that
\[
Q((-\infty, x]) = P((-\infty, x]) \qquad \forall x \in \mathbb{R},
\]
then $P(B) = Q(B)$ for all $B \in \mathcal{B}(\mathbb{R})$.
\end{thmbox}

\noindent\textit{Proof} We can apply the uniqueness of measure extension theorem to the algebra
\[
\mathcal{B}_0 = \{\text{finite disjoint union of sets of the form } (a, b]\}.
\]
When $a < b$, the probability of $(a, b]$ with respect to $P$ is given by
\[
P((a, b]) = P((-\infty, b]) - P((-\infty, a]) = F(b) - F(a).
\]
By extending this to finite disjoint union, we obtain a pre-measure, denoted by $\mu_0$, on the sets in $\mathcal{B}_0$. In fact, if $A = \uplus_{i=1}^n (a_i, b_i]$, with $a_1 < b_1 \le a_2 < b_2 \le \dots \le a_n < b_n$, then
\[
\mu_0(A) = P(A) = \sum_{i=1}^n [F(b_i) - F(a_i)].
\]
Suppose $Q$ is a probability measure identical to $P$ when restricted to sets in the form $(-\infty, a]$. Then for any $A \in \mathcal{B}_0$, we have
\[
Q(A) = \sum_{i=1}^n [F(b_i) - F(a_i)] = P(A).
\]
Since $\mathcal{B}(\mathbb{R})$ is the $\sigma$-field generated by $\mathcal{B}_0$, the extensions $Q$ and $P$ of $\mu_0$ must be the same probability measure by Theorem 3.8. \hfill $\blacksquare$